\documentclass{article}
% SYNTHETIC toy paper for the ScholarForge demo. It has no real content and no real results.
% Every numeric macro below must resolve to a value in results/*.json through macros_map.json.
% Two problems are planted on purpose: a stale macro (accBase) and a hard-coded number in Section 3.

\newcommand{\accMain}{81.4}     % results/main.json: metrics.accuracy x 100
\newcommand{\accBase}{73.9}     % results/baseline.json: metrics.accuracy x 100 (stale on purpose)
\newcommand{\gainPP}{7.2}       % results/main.json: metrics.gain_pp
\newcommand{\seStd}{1.2}        % results/main.json: metrics.stderr_pp
\newcommand{\nSeeds}{5}         % results/main.json: run.n_seeds
\newcommand{\nExamples}{160}    % results/main.json: run.n_examples (defined, never used: soft advisory)

\begin{document}

\title{A Synthetic Paper About Nothing}
\author{Anonymous Authors}
\maketitle

\begin{abstract}
This is a synthetic abstract. The toy method reaches \accMain\% accuracy against
\accBase\% for the toy baseline, a gain of \gainPP{} points over \nSeeds{} seeds.
\end{abstract}

\section{Introduction}
It is worth noting that this paper is a placeholder. We delve into a toy task
and leverage a toy model to unlock a toy result.

\section{Method}
The method is a stand-in. It has no equations and no code.
Crucially, it exists only to give the checkers something to read.

\section{Results}
Table~\ref{tab:main} lists the two systems. The mean difference is \gainPP{} points
with a standard error of \seStd{} points.
Our method also reaches 79.5\% on the held-out split. % planted: hard-coded number

\begin{table}[h]
\centering
\begin{tabular}{lr}
System & Accuracy (\%) \\
Baseline & \accBase \\
Toy method & \accMain \\
\end{tabular}
\caption{Synthetic accuracy of two systems.}
\label{tab:main}
\end{table}

\section{Conclusion}
The toy method wins --- in a toy world.

\end{document}
